\chapter{Experimental Setup and Reproducibility Protocol}
\label{chap:experimental-setup}

\section{Benchmark Environment and Specifications}
To ensure reproducibility, all experiments were conducted within an isolated virtual environment with locked dependency versions:
\begin{itemize}
    \item \textbf{Operating System:} Microsoft Windows (AMD64)
    \item \textbf{Python Runtime:} Python 3.12.13
    \item \textbf{Qiskit SDK:} Qiskit version 2.5.2
    \item \textbf{Simulation Backend:} Qiskit Aer 0.17.2 with reference primitives (\texttt{StatevectorEstimator} and \texttt{StatevectorSampler})
    \item \textbf{Scientific Stack:} NumPy 2.2.3, SciPy 1.18.1, Matplotlib 3.11.2, Rustworkx 0.18.1
    \item \textbf{Git Provenance:} Commit \texttt{73c834927348f9ee606d04a06e0be5582d26e8aa}
    \item \textbf{Artifact Path:} \path{outputs/runs/20261011_103427_qaoa_n5_seed42/experiment_metrics.json}
\end{itemize}

\section{Calibrated 5-Asset Renewable Portfolio Instance}
The benchmark instance models a realistic municipal clean energy capital allocation dilemma. The candidate catalog consists of $N = 5$ renewable assets described in Table \ref{tab:asset_catalog}:
\begin{itemize}
    \item Asset 0: Utility Solar Farm A ($r_0 = 0.85$, $\Sigma_{00} = 0.165$)
    \item Asset 1: Offshore Wind Array B ($r_1 = 0.92$, $\Sigma_{11} = 0.174$)
    \item Asset 2: Lithium-Ion BESS C ($r_2 = 0.65$, $\Sigma_{22} = 0.188$)
    \item Asset 3: Biomass Cogen Plant D ($r_3 = 0.70$, $\Sigma_{33} = 0.142$)
    \item Asset 4: Geothermal Plant E ($r_4 = 0.88$, $\Sigma_{44} = 0.159$)
\end{itemize}

\begin{table}[htbp]
\centering
\small
\caption{Calibrated 5-asset renewable infrastructure candidate catalog.}
\label{tab:asset_catalog}
\resizebox{\textwidth}{!}{%
\begin{tabular}{cllccc}
\toprule
\textbf{Index ($i$)} & \textbf{Asset Name} & \textbf{Technology Category} & \textbf{Return ($r_i$)} & \textbf{Carbon Offset (kt/yr)} & \textbf{Variance ($\Sigma_{ii}$)} \\
\midrule
0 & Utility Solar Farm A   & Solar PV         & 0.85 & 95.0  & 0.165 \\
1 & Offshore Wind Array B  & Offshore Wind    & 0.92 & 120.0 & 0.174 \\
2 & Lithium-Ion BESS C     & Battery Storage  & 0.65 & 60.0  & 0.188 \\
3 & Biomass Cogen Plant D  & Biomass          & 0.70 & 75.0  & 0.142 \\
4 & Geothermal Plant E     & Geothermal       & 0.88 & 110.0 & 0.159 \\
\bottomrule
\end{tabular}%
}
\end{table}


The inter-asset covariance matrix $\Sigma$ was constructed via a positive semi-definite latent factor model with fixed random seed 42, reflecting realistic weather correlation: solar and wind exhibit moderate negative diurnal covariance, while battery storage provides low covariance across all assets.

The instance parameters are configured as:
\begin{itemize}
    \item Total Assets: $N = 5$
    \item Capital Budget: $B = 2$ projects (feasible subspace size $|\Omega_{\text{feasible}}| = \binom{5}{2} = 10$)
    \item Risk-Aversion Factor: $q = 0.5$
    \item Penalty Multiplier: $\lambda = 2.5$
    \item Measurement Shots: $S = 4096$
    \item Reproducibility Seed: 42
\end{itemize}

\section{Experimental Design and Parameter Sweeps}
Two controlled experimental series were executed:
\begin{enumerate}
    \item \textbf{Circuit Depth Scaling ($p$-Sweep):}
    QAOA was evaluated across depths $p \in \{1, 2, 3\}$. For each depth, the number of variational parameters scales as $2p \in \{2, 4, 6\}$. We tracked convergence trajectory, optimal utility, approximation ratio, feasibility yield, and total execution time.
    
    \item \textbf{Penalty Weight Sensitivity ($\lambda$-Sweep):}
    Holding depth fixed at $p = 1$, the penalty multiplier was varied over seven values:
    \begin{equation}
    \lambda \in \{0.5, 1.0, 1.5, 2.0, 2.5, 3.5, 5.0\}
    \end{equation}
    This sweep evaluates the transition between under-penalization (infeasible constraint leakage) and over-penalization (local minimum trapping).
\end{enumerate}

\section{Verification of Hamiltonian Ground-State Equivalence}
Prior to executing QAOA, the compiled \texttt{SparsePauliOp} was formally validated against the classical cost function across all $2^5 = 32$ computational basis states. The diagnostic confirmed that the maximum energy deviation $|\langle x | \hat{H}_C | x \rangle - C(x)| < 10^{-12}$, confirming analytical equivalence.
